\begin{abstract}
An intersection cut depends on both the corner relaxation and the convex
set used to exclude its vertex. We study how this choice affects the
objective bound for quadratic constraints. We first give a self-contained
account of the unrestricted benchmark: with positive reduced costs, the
supremum over all convex free sets equals the nonlinear corner bound.
We distinguish this equality from attainment and give a sufficient
transversality condition for attainment. For bilinear constraints, we
optimize a three-parameter family obtained from determinant-preserving
transformations, and its upward completions. A scaled measure of vertex
depth gives a positive approximation guarantee whose order is sharp.
For a fixed quadratic representation, the point rule used by SCIP admits
a sharp guarantee that also depends on ray conditioning. Exact rational
examples show that the transformed families can miss the corner bound,
even at a unique minimizing point reached on one ray, and that taking every cut from a
family need not remove the gap. We give a contact criterion that separates
supremum equality from attainment, and extend the orbit analysis to
homogeneous $2\times2$ minors. For successive cutting rounds, a uniform
depth condition ensures convergence on a compact relaxation; an explicit
sequence of maximal bilinear-free sets shows why maximality and
uniformly pointed basis cones alone are insufficient. Archived computational
results compare corner-bound selection, reoptimized LP bounds, and solver
performance. They show that improving the first criterion does not
consistently improve the others.
\end{abstract}

\noindent\textbf{Keywords:} intersection cuts; quadratic optimization;
bilinear constraints; free sets; corner relaxations; cutting-plane convergence.
